\section{Fazit}
Zusammenfassend kann gesagt werden, dass der Versuch erfolgreich die Auswirkungen unterschiedlicher Montagearten auf die Wärmeableitung eines Leistungshalbleiters demonstriert hat. Die Messungen zeigten deutlich, dass der Einsatz einer Isolierscheibe den Wärmewiderstand zwischen Gehäuse und Kühlkörper signifikant erhöht, was zu höheren Sperrschichttemperaturen und längeren Zeitkonstanten führt. 
Die experimentell ermittelten Wärmewiderstände wichen jedoch teilweise sehr stark von den theoretisch berechneten Werten ab, was auf verschiedene Fehlerquellen zurückgeführt werden kann, wie z.B. unzureichende Befestigung des Kühlkörpers oder Messungenauigkeiten.